%%% chapters/00_abstract.tex — 요약(Abstract)
\begin{abstract}
 \noindent 본 연구는 왕복 번역(Round-Trip Translation, RTT) 기법과 대형 언어 모델(LLM)을 활용하여 프로그래밍 언어 간의 의미적 거리를 정량적으로 측정하는 방법을 제시한다. 기준이 되는 \cpp{} 코드를 C, Java, Python으로 번역하고 이를 다시 원본 언어로 역번역하는 순환 과정을 통해 의미적 거리를 정량적으로 측정했다. \code{qwen2.5-coder:7b}와 \code{gpt-5.4}를 대상으로 총 30건의 실험을 수행했다. 두 모델은 뚜렷한 차이를 드러냈다. 코드를 전문적으로 생성하는 소형 모델인 Qwen은 원본의 구문적 구조를 최대한 보존하려는 보수적 전략을 취하며 언어별 편차 없이 평균 53.3\%의 균형 잡힌 성공률을 기록했다. 반면, GPT-5.4는 cpp\rarr{}c} 변환 과정에서 의미적 등가성을 100\% 완벽하게 유지하면서도, 원본의 구조에 얽매이지 않고 대상 언어의 패러다임에 맞게 코드를 전면 재구성하는 창조적 재해석의 경향을 보였다. 본 연구는 정점(Fixed-point) 도달 여부, 반복 횟수, 잔차 유사도(Residual Similarity), 추상 구문 트리(AST) 기반의 구조적 거리, 그리고 MOSS 유사도 등 다차원적인 지표를 통합하여 분석했다. RTT 프레임워크는 향후 소프트웨어 엔지니어링에서 이기종 언어 간의 코드 마이그레이션 난이도를 사전에 예측하고, AI 기반 코드 변환 도구의 신뢰성을 검증하는 지침과 도구가 될 것으로 기대한다.

  \vspace{0.5em}
  \noindent\textbf{키워드:} 왕복 번역(Round-Trip Translation), 프로그래밍 언어 거리, 대규모 언어 모델(LLM), 코드 변환, AST 편집 거리, 잔차 유사도
\end{abstract}
